\documentclass[UTF8,fontset=none,11pt]{ctexart}
\usepackage[a4paper,margin=2.05cm]{geometry}
\usepackage{amsmath,amssymb}
\usepackage{enumitem}
\usepackage[table]{xcolor}
\usepackage{booktabs}
\usepackage{array}
\usepackage{longtable}
\usepackage{tikz}
\usepackage{graphicx}
\usepackage{hyperref}
\usetikzlibrary{positioning,arrows.meta,shapes.geometric,fit,backgrounds,calc}

\setCJKmainfont[AutoFakeSlant=0.18]{PingFang SC}
\setCJKsansfont[AutoFakeSlant=0.18]{PingFang SC}
\setCJKmonofont[AutoFakeSlant=0.18]{PingFang SC}

\hypersetup{
  colorlinks=true,
  linkcolor=blue!55!black,
  urlcolor=blue!55!black,
  citecolor=blue!55!black,
  pdftitle={Code Optimization 中文学习笔记},
  pdfauthor={Trae AI}
}

\setlength{\parindent}{0pt}
\setlength{\parskip}{0.55em}
\setlist[itemize]{leftmargin=1.8em,itemsep=0.22em,topsep=0.25em}
\setlist[enumerate]{leftmargin=2.0em,itemsep=0.28em,topsep=0.3em}
\renewcommand{\arraystretch}{1.18}
\setlength{\emergencystretch}{3em}
\sloppy

\definecolor{defrule}{RGB}{25,92,140}
\definecolor{noterule}{RGB}{120,80,15}
\definecolor{warnrule}{RGB}{160,50,45}
\definecolor{softblue}{RGB}{235,246,255}
\definecolor{softyellow}{RGB}{255,249,230}
\definecolor{graytext}{RGB}{95,95,95}
\definecolor{GRAYTEXT}{RGB}{95,95,95}
\newcolumntype{L}[1]{>{\raggedright\arraybackslash}p{#1}}

\newenvironment{defbox}[1][]{%
  \par\medskip\noindent{\color{defrule}\rule{\linewidth}{1.1pt}}\par\nopagebreak
  \noindent{\bfseries\color{defrule}#1}\par\nopagebreak\smallskip}%
  {\par\nopagebreak\smallskip\noindent{\color{defrule}\rule{\linewidth}{0.4pt}}\par\medskip}
\newenvironment{notebox}[1][]{%
  \par\medskip\noindent{\color{noterule}\rule{\linewidth}{1.1pt}}\par\nopagebreak
  \noindent{\bfseries\color{noterule}#1}\par\nopagebreak\smallskip}%
  {\par\nopagebreak\smallskip\noindent{\color{noterule}\rule{\linewidth}{0.4pt}}\par\medskip}
\newenvironment{warnbox}[1][]{%
  \par\medskip\noindent{\color{warnrule}\rule{\linewidth}{1.1pt}}\par\nopagebreak
  \noindent{\bfseries\color{warnrule}#1}\par\nopagebreak\smallskip}%
  {\par\nopagebreak\smallskip\noindent{\color{warnrule}\rule{\linewidth}{0.4pt}}\par\medskip}

\newcommand{\pg}[1]{\texorpdfstring{{\small\color{graytext}\,[p#1]}}{ [p#1]}}
\newcommand{\IR}{\ensuremath{\mathrm{IR}}}
\newcommand{\CFG}{\ensuremath{\mathrm{CFG}}}
\newcommand{\DAG}{\ensuremath{\mathrm{DAG}}}
\newcommand{\CSE}{\ensuremath{\mathrm{CSE}}}
\newcommand{\DCE}{\ensuremath{\mathrm{DCE}}}
\newcommand{\GCP}{\ensuremath{\mathrm{GCP}}}
\newcommand{\TAC}{\ensuremath{\mathrm{TAC}}}
\newcommand{\eps}{\varepsilon}

\title{\bfseries Lecture 9: Intermediate Code Optimization\\中文学习笔记}
\author{基于课件 \texttt{8.Code\_Optimization.pdf} 整理}
\date{\today}

\begin{document}
\maketitle
\tableofcontents
\newpage

\section*{阅读导引}
\addcontentsline{toc}{section}{阅读导引}

本讲讲中间代码优化。前一讲已经把带标注 AST 翻译成 IR，本讲关心如何在保持程序语义等价的前提下改写 IR，使程序更快、更小、更省内存或能耗更低。主线是：先理解优化目标和分类，再构造基本块与控制流图，随后在基本块内做局部优化，在 CFG 上做全局优化，最后联系 LLVM Pass 和优化等级。

\begin{notebox}[本讲主线]
代码优化不是随意“改代码”，而是在证明语义等价的前提下做变换。局部优化依赖基本块和 DAG，全局优化依赖 CFG 和数据流分析。编译器宁可保守地放弃一次优化，也不能做一次可能改变结果的优化。
\end{notebox}

\section{编译阶段与优化入口 \pg{1--3}}

代码优化位于中间代码生成之后、目标代码生成之前：
\[
\text{Semantic Analyzer}\to \text{Intermediate Code Generation}\to \text{Code Optimizer}\to \text{Code Generator}.
\]

优化器的输入通常是中间表示 IR 和符号表，输出仍是 IR，只是这份 IR 在执行时间、空间、能耗或代码大小上更好。

课件第 3 页的 while 示例：

\begin{verbatim}
while (y < z) {
  int x = a + b;
  y += x;
}
\end{verbatim}

中间代码中 \(a+b\) 被计算后放入临时变量 \(t1\)，并移到循环外复用：

\begin{verbatim}
t1 = a + b
goto L1
L2:
t2 := y + t1
y := t2
L1:
if y < z goto L2
\end{verbatim}

这体现了循环不变代码外提的思想：若某个表达式在循环中每次结果都相同，就可以在循环前计算一次，避免重复计算。

\section{优化目标、发生位置与基本原则 \pg{4--6}}

\subsection{优化器做什么 \pg{4}}

课件把优化器分成三类活动：

\begin{itemize}
  \item 控制流分析（control-flow analysis）：弄清程序可能按哪些路径执行。
  \item 数据流分析（data-flow analysis）：弄清变量值、定义、使用等信息如何沿路径传播。
  \item 代码变换（code transformation）：在分析信息支持下修改 IR。
\end{itemize}

\begin{defbox}[代码优化]
代码优化是在保持程序可观察行为不变的前提下，对程序表示做变换，以提升效率和性能。目标可包括减少执行时间、内存使用、能耗、二进制大小或计算复杂度。
\end{defbox}

\subsection{优化发生在哪里 \pg{5}}

\begin{longtable}{L{0.21\linewidth}L{0.36\linewidth}L{0.33\linewidth}}
\toprule
位置 & 例子 & 特点 \\
\midrule
人工源码层面 & 选择更合适算法和数据结构，写出编译器容易优化的代码 & 程序员负责，收益可能很大，但依赖经验 \\
编译器 IR 层面 & 生成更高效 TAC/SSA，做 CSE、DCE、常量传播等 & 本讲重点，机器无关，便于复用 \\
目标代码生成层面 & 选择指令、分配寄存器、调度指令 & 依赖目标机器 ISA 和硬件特性 \\
汇编器/链接器层面 & 链接时优化、重排代码布局 & 编译后期，能利用跨文件或二进制级信息 \\
\bottomrule
\end{longtable}

\subsection{优化原则 \pg{6}}

优化希望改善以下一个或多个指标：

\begin{itemize}
  \item 执行时间。
  \item 内存使用。
  \item 能源消耗。
  \item 二进制可执行文件大小。
\end{itemize}

但优化必须遵守三个原则：

\begin{itemize}
  \item 等价性：优化前后程序结果不能改变。这是最重要的正确性要求。
  \item 高效性：优化后目标代码应更快或占用更少存储空间。
  \item 成本效益：优化本身不能太昂贵，否则编译时间过长可能得不偿失。
\end{itemize}

\section{优化分类：布局相关与代码相关 \pg{7--14}}

\subsection{优化本质是程序变换 \pg{7--9}}

编译器优化本质上是一类程序变换，包括：

\begin{itemize}
  \item 删除：删掉不影响结果的代码。
  \item 添加：插入能提升性能的辅助指令。
  \item 移动：调整指令顺序或代码布局。
  \item 修改：把昂贵操作替换成便宜操作。
\end{itemize}

课件把优化分成两大类：

\begin{longtable}{L{0.20\linewidth}L{0.35\linewidth}L{0.35\linewidth}}
\toprule
类别 & 目标 & 例子 \\
\midrule
Layout-related transformations & 优化代码和数据在内存中的放置，最大化空间局部性 & 函数布局、字段顺序调整、AOS 转 SOA \\
Code-related transformations & 优化生成什么代码，尽量执行更少、更便宜的指令 & CSE、DCE、常量折叠、常量传播、强度削减 \\
\bottomrule
\end{longtable}

\begin{defbox}[空间局部性]
空间局部性指：访问某个内存位置后，附近位置很可能很快也被访问。缓存一次会装入一整条 cache line，因此若相关代码或数据靠得近，就更容易命中缓存。
\end{defbox}

\subsection{代码布局示例 \pg{10}}

课件假设数据 cache 一条 line 有 \(N\) 个 word，每个函数大小为 \(N/2\) word，访问序列为：
\[
g,\ f,\ h,\ f,\ h,\ f,\ h.
\]

若代码布局为 \(f,g,h\)，访问时可能产生 6 次 cache miss；若布局为 \(f,h,g\)，由于 \(f\) 和 \(h\) 放在同一 cache line 中，访问序列反复在 \(f/h\) 间切换时可复用缓存，miss 可降到 2 次。

通俗理解：把经常相邻执行的函数放近，像把经常一起用的书放在同一个书架格子里，拿一次就能顺带拿到另一部分。

\subsection{数据布局示例 \pg{11--12}}

示例 1：调整结构体字段顺序。

\begin{verbatim}
struct S {
  int a;
  int b[10];
  int c;
} obj[100];

for (...) {
  ... = obj[i].a + obj[i].c;
}
\end{verbatim}

若把字段顺序改为：

\begin{verbatim}
struct S {
  int a;
  int c;
  int b[10];
} obj[100];
\end{verbatim}

那么循环中一起访问的 \(a\) 和 \(c\) 更可能落在同一 cache line 中，访问 \(c\) 更可能命中缓存。

示例 2：AOS 转 SOA。

\begin{verbatim}
// AOS: array of structs
struct S { int a; int b; } objs[100];

// SOA: struct of arrays
struct S { int a[100]; int b[100]; } objs;
\end{verbatim}

如果两个循环分别只访问所有 \(a\) 或所有 \(b\)，SOA 让同字段连续排列，空间局部性更好。

\section{代码相关优化的四类变换 \pg{15}}

\begin{longtable}{L{0.18\linewidth}L{0.36\linewidth}L{0.36\linewidth}}
\toprule
变换 & 含义 & 示例 \\
\midrule
Modifying & 修改指令形式，用更便宜指令替代昂贵指令 & \(A=2*a\) 改为 \(A=a<<1\) \\
Deleting & 删除不影响结果的指令 & \(A=2; A=y;\) 中前一条赋值可删 \\
Moving & 调整独立指令顺序，改善并行性或隐藏延迟 & \(A=x/y; B=A+1; C=y;\) 改为 \(A=x/y; C=y; B=A+1;\) \\
Inserting & 插入有助于性能的指令 & 链表遍历中提前 \texttt{prefetch(p->next)} \\
\bottomrule
\end{longtable}

这些变换对应常见优化：

\begin{itemize}
  \item 强度削减（strength reduction）是 modifying。
  \item 死代码删除（dead code elimination）是 deleting。
  \item 代码调度（code scheduling）是 moving。
  \item 数据预取（data prefetching）是 inserting。
\end{itemize}

\section{控制流分析、基本块与 CFG \pg{16--22}}

\subsection{控制流分析 \pg{16}}

\begin{defbox}[Control-Flow Analysis]
控制流分析是静态分析程序执行路径的方法。其结果通常表示为控制流图 CFG，用于后续全局分析和优化。
\end{defbox}

命令式语言中控制流通常由 if、while、for、goto、return 等结构显式决定。要构造 CFG，第一步是把代码划分为基本块。

\subsection{基本块 \pg{17}}

\begin{defbox}[Basic Block]
基本块是直线代码序列：除第一条指令外没有其他标签或跳转入口；除最后一条指令外没有跳转。也就是说，一个基本块只有一个入口和一个出口。
\end{defbox}

基本块的性质：

\begin{itemize}
  \item 只能从块开头进入。
  \item 只能从块末尾跳出。
  \item 块内指令要么全部执行，要么都不执行。
  \item 基本块不能再被进一步切分为更小的控制流单元。
  \item CFG 中的节点就是基本块。
\end{itemize}

局部优化只在单个基本块内进行；全局优化跨基本块，因此需要 CFG。

\subsection{控制流图 \pg{18}}

\begin{defbox}[Control Flow Graph]
控制流图是有向图：每个节点代表一个基本块，每条有向边代表执行流可从一个基本块转移到另一个基本块。
\end{defbox}

边可能来自：

\begin{itemize}
  \item 一个基本块结束后自然落入下一个基本块。
  \item 条件分支导致控制流分叉。
  \item 多条路径汇合到同一基本块。
\end{itemize}

CFG 广泛用于表示单个函数，并支撑全局程序分析和优化。

\subsection{构造 CFG \pg{19--22}}

第一步：找 leader。满足以下任一条件的指令是 leader：

\begin{enumerate}
  \item 程序第一条指令。
  \item 条件或无条件跳转的目标指令。
  \item 紧跟在条件或无条件跳转之后的指令。
\end{enumerate}

从一个 leader 开始，到下一个 leader 之前的所有指令构成一个基本块。

第二步：添加基本块之间的边。若 \(B_1\) 和 \(B_2\) 满足以下情况，则添加边：

\begin{itemize}
  \item \(B_2\) 紧跟 \(B_1\)，且 \(B_1\) 可能 fall through 到 \(B_2\)。条件跳转的 false 分支和非跳转结束都可能 fall through；无条件跳转不会 fall through。
  \item \(B_1\) 以跳转到 \(B_2\) 结束，即使 \(B_2\) 不紧跟 \(B_1\)。
\end{itemize}

课件奇偶判断例子：

\begin{verbatim}
0: t0 = read_num
1: if t0 mod 2 == 0
2: print t0 + " is even."
3: goto 5
4: print t0 + " is odd."
5: end program
\end{verbatim}

leaders 为 \(0,2,4,5\)：0 是第一条，2 和 5 是跳转目标，4 是紧跟跳转后的指令。

\section{局部优化与 DAG \pg{23--33}}

\subsection{局部优化与全局优化 \pg{23}}

\begin{defbox}[局部优化]
局部优化只在一个基本块内部进行。由于基本块内没有分支，块内指令要么按顺序执行一次，要么完全不执行，因此局部优化最简单。
\end{defbox}

\begin{defbox}[全局优化]
全局优化跨基本块进行，作用域可包含 if、while、for 等控制流结构。这里的 global 通常指函数内跨基本块，不一定跨整个程序。
\end{defbox}

局部优化例子：

\begin{itemize}
  \item 常量折叠：\texttt{const int a = 2 + 3;} 可变为 \texttt{a=5}，再用于后续表达式。
  \item 公共子表达式消除：同一个基本块内重复计算的表达式可复用前一次结果。
\end{itemize}

\subsection{公共子表达式消除与死代码删除 \pg{25}}

\begin{defbox}[Common Subexpression Elimination]
公共子表达式消除把多次计算同一值的表达式改为只计算一次，后续使用第一次结果。
\end{defbox}

课件例子：

\begin{verbatim}
y = x + z;
y = x * x + (x / 3);
z = x * x + y;
\end{verbatim}

拆成 TAC 后，\(x*x\) 被计算两次：

\begin{verbatim}
t1 = x * x
t2 = x / 3
y = t1 + t2
t3 = x * x
z = t3 + y
\end{verbatim}

可优化为：

\begin{verbatim}
t1 = x * x
t2 = x / 3
y = t1 + t2
z = t1 + y
\end{verbatim}

\begin{defbox}[Dead Code Elimination]
若一条指令的结果之后从未被使用，则该指令是死代码，可从指令流中删除。
\end{defbox}

死代码删除常和其他优化连用：CSE、常量传播之后常会产生新的死代码。

\subsection{基本块 DAG \pg{26--29}}

\begin{defbox}[DAG of Basic Block]
基本块 DAG 用有向无环图表示基本块内部的值流。叶子节点代表变量初始值，中间节点代表运算，节点可附着一个或多个变量名。它能帮助识别公共子表达式和死代码。
\end{defbox}

构造算法：

\begin{enumerate}
  \item 为基本块中出现的变量初始值创建叶子节点。
  \item 对每条语句创建与该语句关联的节点 \(N\)。
  \item \(N\) 的子节点是该语句操作数在此语句之前的最近定义。
  \item 用语句中的运算符标记 \(N\)。
  \item 某些节点被指定为输出节点，表示块出口仍需保留的值。
\end{enumerate}

课件基本块：

\begin{verbatim}
(1) a = b + c
(2) b = a - d
(3) c = b + c
(4) d = a - d
\end{verbatim}

可用 DAG 发现：

\begin{itemize}
  \item 若 \(b\) 在块出口不活跃，则语句 (2) 的结果可不用保留。
  \item 若 \(b\) 和 \(d\) 都在出口活跃，(2) 与 (4) 都计算 \(a-d\)，可只计算一次，再复制给另一个变量。
  \item 若只有 \(a\) 在出口活跃，则与 \(b,c,d\) 最终值相关的节点都可视为死代码删除。
\end{itemize}

\subsection{代数恒等式与强度削减 \pg{30--31}}

普通 CSE 只识别“语法上明显相同”的表达式；但有些表达式虽然写法不同，数学上等价。课件例子：

\begin{verbatim}
(1) a = b + c
(2) b = b - d
(3) c = c + d
(4) e = b + c
\end{verbatim}

第 (1) 和第 (4) 的 \(b+c\) 在数学上相等，因为 \((b-d)+(c+d)=b+c\)。普通 DAG 可能错过，需要代数恒等式帮助。

常见代数恒等式：

\[
a*1\equiv a,\qquad a*0\equiv 0,\qquad b\ \&\ true\equiv b.
\]

重结合和交换律：
\[
(a+b)+c\equiv a+(b+c),\qquad a+b\equiv b+a.
\]

\begin{defbox}[Strength Reduction]
强度削减是把昂贵操作替换成便宜操作，例如用移位、加法替代乘法或除法。
\end{defbox}

例子：

\begin{verbatim}
x = y / 8  ->  x = y >> 3
y = y * 8  ->  y = y << 3
x^2        ->  x * x
2 * x      ->  x + x
\end{verbatim}

\subsection{常量折叠与常量传播 \pg{32--33}}

\begin{defbox}[Constant Folding]
常量折叠是在编译期直接计算常量表达式，避免运行时重复计算。
\end{defbox}

例子：

\begin{verbatim}
#define LEN 100
x = 2 * LEN;
if (LEN < 0) print("error");
\end{verbatim}

可折叠为：

\begin{verbatim}
x = 200;
if (false) print("error");
\end{verbatim}

随后 DCE 可删除永远不会执行的 if。

\begin{defbox}[Constant Propagation]
常量传播是在编译期把已知常量值替换到变量使用处。
\end{defbox}

局部例子：

\begin{verbatim}
x = 4
y = x * 3
\end{verbatim}

可传播并折叠为：

\begin{verbatim}
x = 4
y = 12
\end{verbatim}

常量传播有局部版本 LCP 和全局版本 GCP。GCP 更强，因为它跨基本块，但也更复杂，需要 CFG 和数据流分析。

\section{全局优化、保守性与数据流分析 \pg{34--41}}

\subsection{全局优化必须对所有路径正确 \pg{35--36}}

全局优化跨越控制流，因此必须考虑所有可能路径。

可优化例子：

\begin{verbatim}
X = 3;
if (B > 0)
  Y = Z + W;
Y = 0;
A = 2 * X;
\end{verbatim}

从 \(X=3\) 到 \(A=2*X\) 的所有路径上，\(X\) 都没有被改写，所以可优化为：

\begin{verbatim}
A = 2 * 3;
\end{verbatim}

不可优化例子：

\begin{verbatim}
X = 3;
if (B > 0)
  Y = Z + W;
X = Y;
Y = 0;
A = 2 * X;
\end{verbatim}

这里 \(X\) 可能被 \(Y\) 改写，后面的 \(X\) 不一定仍为 3，因此不能传播。

\subsection{保守性原则 \pg{37}}

\begin{defbox}[Conservative Analysis]
保守分析要求：只有当编译器能证明某性质在某程序点沿所有路径都成立时，才允许基于该性质优化。不确定时必须回答 don't know，并放弃优化。
\end{defbox}

保守性会损失部分优化机会，但保证正确性。编译器优化的底线是不能改变程序语义。

全局常量传播需要同时知道：

\begin{itemize}
  \item 控制流：从哪里可能到哪里。
  \item 数据流：变量值是否在路径上被改写。
\end{itemize}

\subsection{数据流分析 \pg{38}}

\begin{defbox}[Dataflow Analysis]
数据流分析是在程序每个点计算某种属性值的编译器分析框架。属性值沿 CFG 边传播，语句可改变这些值。为了可行且正确，数据流框架通常使用保守近似。
\end{defbox}

以 GCP 为例，程序点处的属性可以表示为：

\[
\{A=?,\ B=10,\ C=?\}.
\]

若有语句：
\begin{verbatim}
A = B + C
\end{verbatim}

且分析知道 \(B=10\)，则可优化为：

\begin{verbatim}
A = 10 + C
\end{verbatim}

\subsection{全局常量传播 \pg{39--41}}

GCP 用以下状态描述变量：

\begin{longtable}{L{0.20\linewidth}L{0.68\linewidth}}
\toprule
状态 & 含义 \\
\midrule
\(x=*\) & 未赋值或默认未知初始状态 \\
\(x=1,x=2,\ldots\) & 已知变量被赋为某个常量 \\
\(x=\#\) & 变量可能来自多个不同值，无法确定为单一常量 \\
\bottomrule
\end{longtable}

算法思想：

\begin{enumerate}
  \item 所有变量初始状态为 \(*\)。
  \item 沿 CFG 反复传播和更新状态。
  \item 若不同路径给出不同常量，则合并为 \(\#\)。
  \item 迭代直到状态不再变化，即达到固定点。
  \item 固定点后，把确定常量替换到使用点。
\end{enumerate}

课件例子中，常量可传播到 \(X+1\) 和 \(2*X\)。GCP 之后，还可以继续做 DCE，删除不再有用的赋值。例如 \(X=2\) 若最终被后续 \(X=3\) 覆盖且无使用，就可删除。

\section{LLVM Pass 与优化等级 \pg{42--45}}

\subsection{LLVM Pass \pg{42--43}}

\begin{defbox}[LLVM Pass]
LLVM 中，优化以 Pass 形式实现。一个 Pass 遍历程序的某个部分，收集信息或执行变换。它接收 LLVM IR，并执行分析和/或变换。
\end{defbox}

Pass 的特点：

\begin{itemize}
  \item 可用 \texttt{opt} 单独运行某个 Pass。
  \item 可在源代码到二进制代码的编译流程中间执行。
  \item 多个 Pass 的顺序由 Pass Manager 管理。
\end{itemize}

这和本讲模型对应：分析 Pass 提供控制流/数据流信息，变换 Pass 用这些信息改写 IR。

\subsection{优化等级 \pg{44--45}}

\begin{longtable}{L{0.14\linewidth}L{0.76\linewidth}}
\toprule
等级 & 含义 \\
\midrule
\texttt{O0} & 不优化。编译最快，生成代码最利于调试。\\
\texttt{O1} & 介于 O0 和 O2 之间，做一部分低成本优化。\\
\texttt{O2} & 中等优化等级，启用大多数常用优化。\\
\texttt{O3} & 类似 O2，但启用更耗时或可能增大代码体积的优化，目标是更快运行。\\
\texttt{Os} & 类似 O2，但额外优化代码大小。\\
\texttt{Oz} & 比 Os 更激进地减小代码体积。\\
\texttt{O4} & 启用链接时优化；Clang 支持 O4，\texttt{opt} 不支持。\\
\bottomrule
\end{longtable}

课件还展示了不同优化等级下 GCC/LLVM benchmark 性能比较，说明优化等级会真实影响性能，但效果也依赖具体程序、编译器和硬件。

\section{总结与阅读 \pg{46--47}}

课件总结：

\begin{itemize}
  \item 布局相关优化目标是最大化空间局部性，包括代码布局、声明顺序和数据结构组织。
  \item 代码相关优化目标是执行最少、最便宜的指令。
  \item 控制流分析需要基本块和 CFG。
  \item 局部优化使用基本块 DAG，可做 CSE、DCE、代数恒等式、常量折叠和常量传播。
  \item 全局优化跨基本块，需要数据流分析，典型例子是全局常量传播。
\end{itemize}

推荐阅读 Dragon Book：

\begin{itemize}
  \item 第 8.4、8.5 节：基于 DAG 的基本块优化。
  \item 第 9.1、9.6.1、9.6.6 节：循环优化基本概念。
  \item 第 9.2.1--9.2.4 节：数据流方程和 reaching definition analysis。
\end{itemize}

\section{DAG 构造复习 \pg{48--50}}

课件末尾复习 DAG 构造。三类语句：

\[
x=y\ op\ z,\qquad x=op\ y,\qquad x=y.
\]

构造规则：

\begin{enumerate}
  \item 若操作数 \(y\) 尚无节点，创建节点 \(y\)。二元操作中若 \(z\) 尚无节点，也创建节点 \(z\)。
  \item 对二元操作，创建运算节点 OP，左孩子为 \(y\) 节点，右孩子为 \(z\) 节点。
  \item 对一元操作，若已有同操作符、同孩子的节点，可复用；否则创建新节点。
  \item 对复制语句 \(x=y\)，让 \(x\) 附着到 \(y\) 的节点。
  \item 每次重新定义 \(x\)，先从旧节点的标识符列表移除 \(x\)，再把 \(x\) 附到新节点。
\end{enumerate}

例子：

\begin{verbatim}
x = A + B
y = x + 3
z = x * y
w = x * y
u = w
a = 4 / 2
b = - a
\end{verbatim}

DAG 可发现 \(z\) 和 \(w\) 都是 \(x*y\)，而 \(u=w\) 只是复制；\(a=4/2\) 可常量折叠为 \(a=2\)，因此 \(b=-a\) 可变为 \(b=-2\)。若只有 \(z\) 和 \(b\) 在出口活跃，则优化后的代码可为：

\begin{verbatim}
x = A + B
y = x + 3
z = x * y
b = -2
\end{verbatim}

\section{练习题与解答}

\subsection*{题 1：第 3 页 while 示例做了什么优化？}

\textbf{解答：} 做了循环不变代码外提和复用。表达式 \(a+b\) 在循环体中每次结果相同，因此先在循环外计算 \(t1=a+b\)，循环中直接使用 \(t1\)，避免每轮重复执行加法。

\subsection*{题 2：优化必须满足哪些原则？}

\textbf{解答：} 至少满足等价性、高效性和成本效益。等价性要求优化后程序结果不变；高效性要求目标代码更快或更省空间；成本效益要求优化本身开销不能过高。

\subsection*{题 3：AOS 转 SOA 为什么能提升空间局部性？}

\textbf{解答：} AOS 中每个对象的多个字段交错排列；若循环只访问某一字段，会不断跳过其他字段。SOA 把同一字段连续放在数组中，访问该字段时更可能连续命中 cache line。

\subsection*{题 4：找出如下代码的 leaders 并划分基本块。}

\begin{verbatim}
0: t0 = read_num
1: if t0 mod 2 == 0 goto 2
2: print even
3: goto 5
4: print odd
5: end
\end{verbatim}

\textbf{解答：} leaders 为 0、2、4、5。0 是第一条指令；2 和 5 是跳转目标；4 是紧跟跳转之后的指令。基本块为 \([0,1]\)、\([2,3]\)、\([4]\)、\([5]\)。

\subsection*{题 5：局部优化和全局优化的区别是什么？}

\textbf{解答：} 局部优化只在一个基本块内进行，不考虑控制流；基本块内指令顺序执行，分析简单。全局优化跨基本块进行，需要考虑 if/while/for 等控制流路径，通常依赖 CFG 和数据流分析。

\subsection*{题 6：对如下代码做公共子表达式消除。}

\begin{verbatim}
t1 = x * x
t2 = x / 3
y = t1 + t2
t3 = x * x
z = t3 + y
\end{verbatim}

\textbf{解答：} \(t1\) 和 \(t3\) 都计算 \(x*x\)，且中间没有改变 \(x\)。可删掉第二次计算并改用 \(t1\)：

\begin{verbatim}
t1 = x * x
t2 = x / 3
y = t1 + t2
z = t1 + y
\end{verbatim}

\subsection*{题 7：常量折叠与常量传播有什么区别？}

\textbf{解答：} 常量折叠是在编译期计算常量表达式，例如 \(2*100\) 直接变为 200。常量传播是把已知常量值替换到变量使用处，例如 \(x=4; y=x*3\) 先变为 \(y=4*3\)，再可折叠为 \(y=12\)。

\subsection*{题 8：为什么全局优化要保守？}

\textbf{解答：} 全局优化跨控制流路径。若某性质只在部分路径成立，基于它做优化会改变程序结果。因此编译器必须证明该性质沿所有路径成立；若不能证明，就回答 don't know 并放弃优化。

\subsection*{题 9：解释 GCP 中 \texorpdfstring{\(x=*\)、\(x=3\)、\(x=\#\)}{x states} 的含义。}

\textbf{解答：} \(x=*\) 表示未赋值或默认初始状态；\(x=3\) 表示当前能确定 \(x\) 为常量 3；\(x=\#\) 表示 \(x\) 可能来自多个不同值，不能确定为单一常量，因此不能替换。

\subsection*{题 10：对如下代码说明是否可把最后的 X 替换为 3。}

\begin{verbatim}
X = 3;
if (B > 0)
  Y = Z + W;
Y = 0;
A = 2 * X;
\end{verbatim}

\textbf{解答：} 可以。所有从 \(X=3\) 到 \(A=2*X\) 的路径上，\(X\) 都没有被重新定义，因此最后的 \(X\) 一定为 3。

\subsection*{题 11：下面代码能否把最后的 X 替换为 3？}

\begin{verbatim}
X = 3;
if (B > 0)
  Y = Z + W;
X = Y;
Y = 0;
A = 2 * X;
\end{verbatim}

\textbf{解答：} 不能。中间有 \(X=Y\)，会重新定义 \(X\)，而 \(Y\) 不一定为 3，所以不能安全传播常量 3。

\subsection*{题 12：对 DAG Revisit 示例做优化，假设只有 z 和 b 在出口活跃。}

原代码：

\begin{verbatim}
x = A + B
y = x + 3
z = x * y
w = x * y
u = w
a = 4 / 2
b = - a
\end{verbatim}

\textbf{解答：} \(w\) 与 \(z\) 是同一公共子表达式 \(x*y\)，\(u=w\) 在出口不活跃可删除；\(a=4/2\) 可折叠为 \(a=2\)，若 \(a\) 不活跃，则直接生成 \(b=-2\)。优化后：

\begin{verbatim}
x = A + B
y = x + 3
z = x * y
b = -2
\end{verbatim}

\section*{复习建议}
\addcontentsline{toc}{section}{复习建议}

复习本讲建议按三层掌握：

\begin{itemize}
  \item 概念层：优化目标、等价性、空间局部性、基本块、CFG、DAG、数据流分析、保守性。
  \item 局部层：能用 DAG 识别 CSE、DCE、常量折叠、常量传播和代数恒等式。
  \item 全局层：能解释为什么全局优化需要 CFG 和数据流分析，能判断某个常量是否沿所有路径保持不变。
\end{itemize}

最容易混淆的是“能看到某条路径上为常量”和“所有路径上都为常量”。全局优化只允许后者。

\end{document}
